% GENERATED FILE. Edit canonical lessons, then run npm run generate:books.
% canonical-source-hash: fnv1a64:2e59179049b4bef2
% canonical-lessons: KA-C293-sasive, KA-C293-kottambari, KA-C293-karibevu, KA-C293-tometo, KA-C293-kyaret

\chapter{Mustard seeds, Coriander, Curry leaves, Tomato, Carrot}
\label{ch:ka-mustard-seeds-coriander-curry-leaves-tomato-carrot}
\hlchaptermodality{293}

\begin{chapteropening}
\textbf{By the end of this chapter:} \emph{I can say mustard seeds, coriander, curry leaves, a tomato and a carrot.}

Complete the last lesson of chapter 293: I can say mustard seeds, coriander, curry leaves, a tomato and a carrot.
\end{chapteropening}

\section[\texorpdfstring{\kn{ಸಾಸಿವೆ} (sāsive)}{ಸಾಸಿವೆ (sāsive)}]{\kn{ಸಾಸಿವೆ} (sāsive) — mustard seeds}

\label{lesson:KA-C293-sasive}



\begin{quote}
Before the new one: say the Kannada for messy, then the Kannada for convenient.
\end{quote}

\subsection*{You'll want to know: \kn{ಸಾಸಿವೆ}}
\textbf{\kn{ಸಾಸಿವೆ}} — \emph{sāsive} — \textquotedblleft{}mustard seeds\textquotedblright{}.

\begin{grammarlens}[title={What you've built}]
One more word for this chapter.
\end{grammarlens}

\subsection*{Your turn}
\begin{itemize}
\raggedright
  \item \emph{Say it:} \emph{sāsive}
  \item \emph{Say it:} \emph{sāsive}, once more
  \item \emph{Say it:} say \emph{sāsive}
  \item \emph{Recall:} say the Kannada for messy, then the Kannada for convenient, then say \emph{sāsive} again
\end{itemize}

\begin{tcolorbox}[breakable,colback=teal!4,colframe=teal!35!black,title={Before you move on}]
What does \kn{ಸಾಸಿವೆ} mean? (\textquotedblleft{}Mustard seeds\textquotedblright{}.) Say it once more.
\end{tcolorbox}
\section[\texorpdfstring{\kn{ಕೊತ್ತಂಬರಿ} (kottambari)}{ಕೊತ್ತಂಬರಿ (kottambari)}]{\kn{ಕೊತ್ತಂಬರಿ} (kottambari) — coriander}

\label{lesson:KA-C293-kottambari}



\begin{quote}
Before the new one: say the Kannada for convenient, then the Kannada for mustard seeds.
\end{quote}

\subsection*{You'll want to know: \kn{ಕೊತ್ತಂಬರಿ}}
\textbf{\kn{ಕೊತ್ತಂಬರಿ}} — \emph{kottambari} — \textquotedblleft{}coriander\textquotedblright{}.

\begin{grammarlens}[title={What you've built}]
One more word for this chapter.
\end{grammarlens}

\subsection*{Your turn}
\begin{itemize}
\raggedright
  \item \emph{Say it:} \emph{kottambari}
  \item \emph{Say it:} \emph{kottambari}, once more
  \item \emph{Say it:} say \emph{kottambari}
  \item \emph{Recall:} say the Kannada for convenient, then the Kannada for mustard seeds, then say \emph{kottambari} again
\end{itemize}

\begin{tcolorbox}[breakable,colback=teal!4,colframe=teal!35!black,title={Before you move on}]
What does \kn{ಕೊತ್ತಂಬರಿ} mean? (\textquotedblleft{}Coriander\textquotedblright{}.) Say it once more.
\end{tcolorbox}
\section[\texorpdfstring{\kn{ಕರಿಬೇವು} (karibēvu)}{ಕರಿಬೇವು (karibēvu)}]{\kn{ಕರಿಬೇವು} (karibēvu) — curry leaves}

\label{lesson:KA-C293-karibevu}



\begin{quote}
Before the new one: say the Kannada for mustard seeds, then the Kannada for coriander.
\end{quote}

\subsection*{You'll want to know: \kn{ಕರಿಬೇವು}}
\textbf{\kn{ಕರಿಬೇವು}} — \emph{karibēvu} — \textquotedblleft{}curry leaves\textquotedblright{}.

\begin{grammarlens}[title={What you've built}]
One more word for this chapter.
\end{grammarlens}

\subsection*{Your turn}
\begin{itemize}
\raggedright
  \item \emph{Say it:} \emph{karibēvu}
  \item \emph{Say it:} \emph{karibēvu}, once more
  \item \emph{Say it:} say \emph{karibēvu}
  \item \emph{Recall:} say the Kannada for mustard seeds, then the Kannada for coriander, then say \emph{karibēvu} again
\end{itemize}

\begin{tcolorbox}[breakable,colback=teal!4,colframe=teal!35!black,title={Before you move on}]
What does \kn{ಕರಿಬೇವು} mean? (\textquotedblleft{}Curry leaves\textquotedblright{}.) Say it once more.
\end{tcolorbox}
\section[\texorpdfstring{\kn{ಟೊಮೇಟೊ} (ṭomēṭo)}{ಟೊಮೇಟೊ (ṭomēṭo)}]{\kn{ಟೊಮೇಟೊ} (ṭomēṭo) — a tomato}

\label{lesson:KA-C293-tometo}



\begin{quote}
Before the new one: say the Kannada for coriander, then the Kannada for curry leaves.
\end{quote}

\subsection*{You'll want to know: \kn{ಟೊಮೇಟೊ}}
\textbf{\kn{ಟೊಮೇಟೊ}} — \emph{ṭomēṭo} — \textquotedblleft{}a tomato\textquotedblright{}.

\begin{grammarlens}[title={What you've built}]
One more word for this chapter.
\end{grammarlens}

\subsection*{Your turn}
\begin{itemize}
\raggedright
  \item \emph{Say it:} \emph{ṭomēṭo}
  \item \emph{Say it:} \emph{ṭomēṭo}, once more
  \item \emph{Say it:} say \emph{ṭomēṭo}
  \item \emph{Recall:} say the Kannada for coriander, then the Kannada for curry leaves, then say \emph{ṭomēṭo} again
\end{itemize}

\begin{tcolorbox}[breakable,colback=teal!4,colframe=teal!35!black,title={Before you move on}]
What does \kn{ಟೊಮೇಟೊ} mean? (\textquotedblleft{}A tomato\textquotedblright{}.) Say it once more.
\end{tcolorbox}
\section[\texorpdfstring{\kn{ಕ್ಯಾರೆಟ್} (kyāreṭ)}{ಕ್ಯಾರೆಟ್ (kyāreṭ)}]{\kn{ಕ್ಯಾರೆಟ್} (kyāreṭ) — a carrot}

\label{lesson:KA-C293-kyaret}



\begin{quote}
Before the new one: say the Kannada for curry leaves, then the Kannada for a tomato.
\end{quote}

\subsection*{You'll want to know: \kn{ಕ್ಯಾರೆಟ್}}
\textbf{\kn{ಕ್ಯಾರೆಟ್}} — \emph{kyāreṭ} — \textquotedblleft{}a carrot\textquotedblright{}.

\begin{grammarlens}[title={What you've built}]
One more word for this chapter.
\end{grammarlens}

\subsection*{Your turn}
\begin{itemize}
\raggedright
  \item \emph{Say it:} \emph{kyāreṭ}
  \item \emph{Say it:} \emph{kyāreṭ}, once more
  \item \emph{Say it:} say \emph{kyāreṭ}
  \item \emph{Recall:} say the Kannada for curry leaves, then the Kannada for a tomato, then say \emph{kyāreṭ} again
\end{itemize}

\begin{tcolorbox}[breakable,colback=teal!4,colframe=teal!35!black,title={Before you move on}]
What does \kn{ಕ್ಯಾರೆಟ್} mean? (\textquotedblleft{}A carrot\textquotedblright{}.) Say it once more.
\end{tcolorbox}
